\section{Bilan d'utilisation et situation finale}

\subsection{Apports de l'approche Python}

L'approche \path{full_soft/} presente des avantages clairs:

\begin{itemize}
  \item elle est directe et lisible pour comprendre le controle embarque;
  \item elle documente l'environnement Raspberry Pi et les interfaces materielles;
  \item elle separe partiellement les interfaces capteurs/actionneurs de l'algorithme;
  \item elle contient un simulateur Python historique utile comme reference;
  \item elle permet de comprendre la logique de conduite avant l'integration ROS.
\end{itemize}

\subsection{Limites pour la Voiture Blue actuelle}

Pour la Voiture Blue actuelle, cette pile semble moins adaptee que la pile APEX. Elle ne fournit pas de graphe ROS, pas de topics standardises, pas de launch files, pas de bridge Gazebo moderne et pas de mecanisme de capture aussi structure que la pile APEX.

Elle reste pertinente pour comprendre l'historique logiciel et certains choix materiels, mais elle ne semble pas etre le chemin recommande pour les essais actuels.

\subsection{Situation finale}

D'apres l'analyse du depot, \path{full_soft/} doit etre considere comme une base historique et pedagogique. Il peut aider a comprendre le style Voiture Jaune et l'ancien pilotage direct sur Raspberry Pi. Pour la Voiture Blue, la situation finale semble etre une migration progressive vers ROS 2/APEX, avec conservation de \path{full_soft/} comme reference et non comme pile principale.

